\documentclass[11pt]{article}
\usepackage{mathblatt}
\usepackage{multicol}


\begin{document}
\pfheftstil
\pfheftkopf{Dreiecke \textperiodcentered{} Lösungen}{}
\begin{pfloesung}{}
\lz{1.}{c (liegt dem rechten Winkel gegenüber)}{}{}
\end{pfloesung}
\begin{pfloesung}{}
\lz{2.}{100; 10; 144; 12}{}{}
\end{pfloesung}
\begin{pfloesung}{}
\lz{3.}{ja; ja; nein}{}{}
\end{pfloesung}
\begin{pfloesung}{}
\lz{4.}{Gegenkathete a; Ankathete b; Hypotenuse c}{}{}
\end{pfloesung}
\begin{pfloesung}{}
\lz{5.}{sin; cos; tan}{}{}
\end{pfloesung}
\begin{pfloesung}{}
\lz{6.}{x = 4; x = 5; x = 12}{}{}
\end{pfloesung}
\begin{pfloesung}{}
\lz{7.}{$\gamma$ = 70°; Seite c}{}{}
\end{pfloesung}
\begin{pfloesung}{}
\lz{8.}{65°; 106°; 53,7°}{}{}
\end{pfloesung}
\begin{pfloesung}{}
\lz{9.}{ja; nein}{}{}
\end{pfloesung}
\begin{pfloesung}{}
\lz{10.}{rechter Winkel bei C}{ja}{}
\end{pfloesung}
\begin{pfloesung}{}
\lz{11.}{Hypotenuse: $41$ cm}{}{}
\end{pfloesung}
\begin{pfloesung}{}
\lz{12.}{$8{,}5$ cm}{}{}
\end{pfloesung}
\begin{pfloesung}{}
\lz{13.}{plus}{Hypotenuse gesucht}{}
\end{pfloesung}
\begin{pfloesung}{}
\lz{14.}{Umkehrung: Wenn ein Tier vier Beine hat, dann ist es ein Hund}{}{}
\end{pfloesung}
\begin{pfloesung}{}
\lz{15.}{$c = 39$ cm}{Gleichung: $c^2$ $\Rightarrow$ $36^2 + 15^2$; Zwischenergebnis: $c^2$ $\Rightarrow$ 1\,296 + 225 = 1\,521; Wurzel: $\sqrt{1\,521}$ = 39 cm}{}
\end{pfloesung}
\begin{pfloesung}{}
\lz{16.}{Katheten $7$ und $9$ am rechten Winkel, $w$ gegenüber dem rechten Winkel}{}{}
\end{pfloesung}
\begin{pfloesung}{}
\lz{17.}{gilt, gilt, gilt nicht, gilt, gilt nicht – $m$ ist die Hypotenuse}{1; 2; 3; 4; 5; 2; 4}{}
\end{pfloesung}
\begin{pfloesung}{}
\lz{18.}{Die Summe der Flächeninhalte der Kathetenquadrate ist gleich dem Flächeninhalt des Hypotenusenquadrates}{a² + b² = c²}{}
\end{pfloesung}
\begin{pfloesung}{}
\lz{19.}{m² + n² = k²}{}{}
\end{pfloesung}
\begin{pfloesung}{}
\lz{20.}{z² = x² + y²}{}{}
\end{pfloesung}
\begin{pfloesung}{}
\lz{21.}{z = √(x² + y²)}{}{}
\end{pfloesung}
\begin{pfloesung}{}
\lz{22.}{z = √(x² $-$ y²) (vierte Option)}{z² = x² $-$ y²}{}
\end{pfloesung}
\begin{pfloesung}{}
\lz{23.}{w = √(u² + v²) (dritte Option)}{w² = u² + v²}{}
\end{pfloesung}
\begin{pfloesung}{}
\lz{24.}{$9$ cm}{Radius}{}
\end{pfloesung}
\begin{pfloesung}{}
\lz{25.}{Katheten: die Höhe und der Überstand}{}{}
\end{pfloesung}
\begin{pfloesung}{}
\lz{26.}{13 cm}{}{}
\end{pfloesung}
\begin{pfloesung}{}
\lz{27.}{$61$ m}{$60^2 + 11^2$ $\Rightarrow$ 3\,721, Weg}{}
\end{pfloesung}
\begin{pfloesung}{}
\lz{28.}{e = √(196² $-$ 10²) km $\approx$ 195,74 km}{}{}
\end{pfloesung}
\begin{pfloesung}{}
\lz{29.}{x = √(2,40² $-$ 2,20²) m $\approx$ 0,96 m}{}{}
\end{pfloesung}
\begin{pfloesung}{}
\lz{30.}{$\sqrt{97} \approx 9{,}8$ cm}{$c^2$ $\Rightarrow$ 97, c}{}
\end{pfloesung}
\begin{pfloesung}{}
\lz{31.}{$h = 63$ cm}{halbe Grundseite: $\tfrac{32}{2}$ $\Rightarrow$ 16 cm; Gleichung umstellen: $h^2$ $\Rightarrow$ $65^2 - 16^2$; Zwischenergebnis: $h^2$ $\Rightarrow$ 4\,225 - 256 = 3\,969; Wurzel: $\sqrt{3\,969}$ = 63 cm}{}
\end{pfloesung}
\begin{pfloesung}{}
\lz{32.}{$53$ m}{Unterschied: $\sqrt{65^2 - 39^2}$ $\Rightarrow$ $\sqrt{2\,704}$ = 52 m; mit der Handhöhe: 52 + 1 $\Rightarrow$ 53}{}
\end{pfloesung}
\begin{pfloesung}{}
\lz{33.}{x = √(5,7² $-$ 0,4²) m $\approx$ 5,69 m}{}{}
\end{pfloesung}
\begin{pfloesung}{}
\lz{34.}{√(30² + 55²) m $\approx$ 62,65 m}{}{}
\end{pfloesung}
\begin{pfloesung}{}
\lz{35.}{Ja}{Diagonale: $\sqrt{111^2 + 62^2}$ $\Rightarrow$ $\sqrt{16\,165}$ $\approx$ 127{,}1 cm; in Zoll: $\tfrac{127{,}1}{2{,}54}$ $\Rightarrow$ $\approx$ 50 Zoll}{}
\end{pfloesung}
\begin{pfloesung}{}
\lz{36.}{AB = 10√13 $\approx$ 36,1 m}{12² + 34² $\Rightarrow$ 1300}{}
\end{pfloesung}
\begin{pfloesung}{}
\lz{37.}{AD = √553 719 $\approx$ 744 m}{Unterschied: 920² $-$ 541²}{}
\end{pfloesung}
\begin{pfloesung}{}
\lz{38.}{BC = √65 $\approx$ 8,06 m}{Unterschied: 9² $-$ 4² $\Rightarrow$ 65}{}
\end{pfloesung}
\begin{pfloesung}{}
\lz{39.}{AD = √(14,1² $-$ 7,4²) $\approx$ 12,0 m}{Unterschied: 198,81 $-$ 54,76 $\Rightarrow$ 144,05}{}
\end{pfloesung}
\begin{pfloesung}{}
\lz{40.}{FA = 3√9159 $\approx$ 287,1 m}{Unterschied: 384² $-$ 255² $\Rightarrow$ 82 431}{}
\end{pfloesung}
\begin{pfloesung}{}
\lz{41.}{Im Becher mindestens 22 cm $-$ 13 cm = 9 cm; h = √(9² $-$ 3²) cm $\approx$ 8,49 cm}{}{}
\end{pfloesung}
\begin{pfloesung}{}
\lz{42.}{√89,96 + 2 $\approx$ 11,5 cm}{Diagonale √(7² + 6,4²) $\Rightarrow$ √89,96 $\approx$ 9,5 cm; Durchmesser: 6,4 cm}{}
\end{pfloesung}
\begin{pfloesung}{}
\lz{43.}{s = √(3² + 4²) = 5 cm}{}{}
\end{pfloesung}
\begin{pfloesung}{}
\lz{44.}{$5$}{x-Unterschied 3, y-Unterschied 4; $PQ^2$ $\Rightarrow$ 25, PQ}{}
\end{pfloesung}
\begin{pfloesung}{}
\lz{45.}{$41$ cm}{halbe Grundkante: 9 cm; $h_s^2$ $\Rightarrow$ $1\,681, h_s$}{}
\end{pfloesung}
\begin{pfloesung}{}
\lz{46.}{$65$ cm}{33 cm; $s^2$ $\Rightarrow$ $33^2 + 56^2$ = 4\,225, s}{}
\end{pfloesung}
\begin{pfloesung}{}
\lz{47.}{$40$ cm}{41 cm; Unterschied: $\sqrt{41^2 - 9^2}$ $\Rightarrow$ $\sqrt{1\,600}$}{}
\end{pfloesung}
\begin{pfloesung}{}
\lz{48.}{$7$ dm}{$Bodendiagonale^2$ $\Rightarrow$ 13; $Raumdiagonale^2$ $\Rightarrow$ $13 + 6^2$ = 49, Raumdiagonale}{}
\end{pfloesung}
\begin{pfloesung}{}
\lz{49.}{$37$ cm}{halbe Grundseite: 12 cm; $Schenkel^2$ $\Rightarrow$ 1\,369, Schenkel}{}
\end{pfloesung}
\begin{pfloesung}{}
\lz{50.}{$60$ cm}{$h^2$ $\Rightarrow$ $61^2 - 11^2$ = 3\,600, h}{}
\end{pfloesung}
\begin{pfloesung}{}
\lz{51.}{$42$ m}{$AD^2$ $\Rightarrow$ $58^2 - 40^2$ = 1\,764, AD}{}
\end{pfloesung}
\begin{pfloesung}{}
\lz{52.}{$12{,}5$ cm}{halbe Diagonalen: 3{,}5 cm und 12 cm; $Seite^2$ $\Rightarrow$ 156{,}25, Seite}{}
\end{pfloesung}
\begin{pfloesung}{}
\lz{53.}{$41$ cm}{Überstand: (50 - 32) : 2 $\Rightarrow$ 9 cm; $Schenkel^2$ $\Rightarrow$ $9^2 + 40^2$ = 1\,681, Schenkel}{}
\end{pfloesung}
\begin{pfloesung}{}
\lz{54.}{Erst die Bodendiagonale aus Länge und Breite}{Quadrate addieren, Wurzel}{}
\end{pfloesung}
\begin{pfloesung}{}
\lz{55.}{$22{,}8$ m}{$\sqrt{23{,}04}$ = 4{,}8 m; gesamt: 18 + 4{,}8 $\Rightarrow$ 22{,}8}{}
\end{pfloesung}
\begin{pfloesung}{}
\lz{56.}{Höhe von C auf AB halbiert $\gamma$: tan($\gamma$/2) = 3 : 8 = 0,375}{}{}
\end{pfloesung}
\begin{pfloesung}{}
\lz{57.}{Radius 3,2 m und Dachhöhe als Katheten, Mantellinie als Hypotenuse $\Rightarrow$ s = √(3,2² + h²)}{}{}
\end{pfloesung}
\begin{pfloesung}{}
\lz{58.}{tan $\angle$FAE = EF : AE = 5,5 : 9 $\Rightarrow$ $\angle$FAE $\approx$ 31,43°}{}{}
\end{pfloesung}
\begin{pfloesung}{}
\lz{59.}{$50$ cm²}{$a^2 + a^2$ $\Rightarrow$ $10^2, 2a^2$ = $100, a^2$ = 50; 7{,}1 cm; 4a $\approx$ 28{,}3 cm; $a^2$}{}
\end{pfloesung}
\begin{pfloesung}{}
\lz{60.}{Ja}{Diagonale des Kofferbodens: $\sqrt{55^2 + 35^2}$ $\Rightarrow$ $\sqrt{4\,250}$ $\approx$ 65{,}2 cm}{}
\end{pfloesung}
\begin{pfloesung}{}
\lz{61.}{x = √(38² + 22²) m $\approx$ 43,91 m}{}{}
\end{pfloesung}
\begin{pfloesung}{}
\lz{62.}{d = √(5² + 3²) cm = √34 cm $\approx$ 5,8 cm}{}{}
\end{pfloesung}
\begin{pfloesung}{}
\lz{63.}{CS = √(11² + 9²) cm = √202 cm $\approx$ 14,21 cm; tan $\angle$SCA = 9 : 11 $\Rightarrow$ $\angle$SCA $\approx$ 39,29°}{}{}
\end{pfloesung}
\begin{pfloesung}{}
\lz{64.}{$\approx$ 345 m (je Seite √(69,3² $-$ 24²) $\approx$ 65,01 m)}{}{}
\end{pfloesung}
\begin{pfloesung}{}
\lz{65.}{25 + √51,87 $\approx$ 32,2 m}{√(8,6² $-$ 4,7²) $\Rightarrow$ √51,87 $\approx$ 7,20 m}{}
\end{pfloesung}
\begin{pfloesung}{}
\lz{66.}{x = 2√7289 $\approx$ 170,8 cm; $\alpha$ = tan$^{-1}$($\tfrac{8}{85}$) $\approx$ 5,4°}{170² + 16² $\Rightarrow$ 29 156; Anteil: 16 : 170 $\Rightarrow$ $\tfrac{8}{85}$ $\approx$ 0,0941}{}
\end{pfloesung}
\begin{pfloesung}{}
\lz{67.}{h\_s = √45 $\approx$ 6,71 m; A $\approx$ 20,1 m²; tan $\alpha$ = 6 / (3√2), $\alpha$ $\approx$ 54,7°; Volumen 4-mal so groß (72 m³ $\to$ 288 m³)}{}{}
\end{pfloesung}
\begin{pfloesung}{}
\lz{68.}{BC = √20 $\approx$ 4,47; $\beta$ = tan$^{-1}$(2) $\approx$ 63,4°}{2, AC = 4; 2² + 4² $\Rightarrow$ 20; 4 : 2 $\Rightarrow$ 2}{}
\end{pfloesung}
\begin{pfloesung}{}
\lz{69.}{im Zähler: mal}{}{}
\end{pfloesung}
\begin{pfloesung}{}
\lz{70.}{sin}{}{}
\end{pfloesung}
\begin{pfloesung}{}
\lz{71.}{r: G, s: A, t: H}{}{}
\end{pfloesung}
\begin{pfloesung}{}
\lz{72.}{sin $\varepsilon$ = AC / DC; tan $\beta$ = AC / AB}{}{}
\end{pfloesung}
\begin{pfloesung}{}
\lz{73.}{x = 14}{}{}
\end{pfloesung}
\begin{pfloesung}{}
\lz{74.}{$\tfrac{e}{f}$}{$\mathrm{sin}\,\beta$}{}
\end{pfloesung}
\begin{pfloesung}{}
\lz{75.}{$\tfrac{q}{p}$}{$\mathrm{sin}\,\beta$}{}
\end{pfloesung}
\begin{pfloesung}{}
\lz{76.}{sin $\beta$ = r/t}{}{}
\end{pfloesung}
\begin{pfloesung}{}
\lz{77.}{sin $\alpha$ = r/t}{}{}
\end{pfloesung}
\begin{pfloesung}{}
\lz{78.}{tan $\alpha$ = r/s}{}{}
\end{pfloesung}
\begin{pfloesung}{}
\lz{79.}{sin $\gamma$ = u/w}{sin = Gegenkathete : Hypotenuse}{}
\end{pfloesung}
\begin{pfloesung}{}
\lz{80.}{$a$ ist jetzt Ankathete, $b$ Gegenkathete}{}{}
\end{pfloesung}
\begin{pfloesung}{}
\lz{81.}{sin $\beta$ = b/a}{}{}
\end{pfloesung}
\begin{pfloesung}{}
\lz{82.}{Winkel: Umkehrtaste}{}{}
\end{pfloesung}
\begin{pfloesung}{}
\lz{83.}{tan $\alpha$ = 2 $\Rightarrow$ $\alpha$ $\approx$ 63,4°}{}{}
\end{pfloesung}
\begin{pfloesung}{}
\lz{84.}{$\approx 18{,}8^\circ$}{$\tfrac{420}{1\,300}, \alpha$}{}
\end{pfloesung}
\begin{pfloesung}{}
\lz{85.}{tan $\alpha$ = 17,5/7,4 $\Rightarrow$ $\alpha$ $\approx$ 67,1°}{}{}
\end{pfloesung}
\begin{pfloesung}{}
\lz{86.}{cos $\alpha$ = 2,5/9 $\Rightarrow$ $\alpha$ $\approx$ 73,9°}{}{}
\end{pfloesung}
\begin{pfloesung}{}
\lz{87.}{tan $\alpha$ = $\tfrac{9}{100}$ $\Rightarrow$ $\alpha$ $\approx$ 5,14° < 7° $\Rightarrow$ die Bahn im Harz ist steiler}{}{}
\end{pfloesung}
\begin{pfloesung}{}
\lz{88.}{tan $\alpha$ = $\tfrac{10}{198}$ $\Rightarrow$ $\alpha$ $\approx$ 2,89°}{}{}
\end{pfloesung}
\begin{pfloesung}{}
\lz{89.}{steiler als $3{,}4^\circ$, nicht zulässig}{$\tfrac{0{,}32}{5}, \alpha$ $\Rightarrow$ $\approx$ $3{,}7^\circ$}{}
\end{pfloesung}
\begin{pfloesung}{}
\lz{90.}{$\approx 36{,}9^\circ$}{Gegenkathete: $3,$ Ankathete $4: \mathrm{tan}\,\beta$ = $\tfrac{3}{4}, \beta$}{}
\end{pfloesung}
\begin{pfloesung}{}
\lz{91.}{$\angle$MSC $\approx$ 35,84° (MC = 6,5 cm; tan $\angle$MSC = 6,5 : 9)}{}{}
\end{pfloesung}
\begin{pfloesung}{}
\lz{92.}{$\angle$CMD $\approx$ 53,13° (DM = 3 cm; cos $\angle$CMD = 3 : 5)}{}{}
\end{pfloesung}
\begin{pfloesung}{}
\lz{93.}{cos $\alpha$ = $\tfrac{4}{9}$ $\Rightarrow$ $\alpha$ $\approx$ 63,6° $\approx$ 64°}{$\tfrac{4}{9}$ $\Rightarrow$ $\approx$ 0,444}{}
\end{pfloesung}
\begin{pfloesung}{}
\lz{94.}{$\alpha$ = sin$^{-1}$($\tfrac{74}{141}$) $\approx$ 31,7°}{Anteil: 7,4 : 14,1 $\Rightarrow$ $\tfrac{74}{141}$ $\approx$ 0,525}{}
\end{pfloesung}
\begin{pfloesung}{}
\lz{95.}{$\beta$1 = cos$^{-1}$($\tfrac{85}{128}$) $\approx$ 48,4°}{cos $\beta$1 = 255 : 384 = $\tfrac{85}{128}$ $\approx$ 0,664}{}
\end{pfloesung}
\begin{pfloesung}{}
\lz{96.}{$\alpha$ = tan$^{-1}$($\tfrac{14}{9}$) $\approx$ 57,3°}{Unterschied: 15 cm $-$ 6 cm $\Rightarrow$ 9 cm; 14 : 9 $\Rightarrow$ 1{,}56}{}
\end{pfloesung}
\begin{pfloesung}{}
\lz{97.}{Teildreieck aus Seite $b$, Höhe und Stück von $a$: $b$ H, Höhe G, Stück von $a$ A}{}{}
\end{pfloesung}
\begin{pfloesung}{}
\lz{98.}{$\approx 4{,}0$ m}{$4{,}2 \cdot \mathrm{sin}\,71^\circ$}{}
\end{pfloesung}
\begin{pfloesung}{}
\lz{99.}{$\approx 4{,}1$ cm}{$\mathrm{sin}\,31^\circ$ $\Rightarrow$ $\tfrac{x}{8},$ x = $8 \cdot \mathrm{sin}\,31^\circ$}{}
\end{pfloesung}
\begin{pfloesung}{}
\lz{100.}{mehr als $5{,}5$ m, die Leiter reicht}{$6 \cdot \mathrm{sin}\,72^\circ$ $\Rightarrow$ $\approx$ 5{,}7 m}{}
\end{pfloesung}
\begin{pfloesung}{}
\lz{101.}{$\approx 6{,}2$ cm}{Teildreieck: Schenkel $H,$ Höhe $G. \mathrm{sin}\,63^\circ$ $\Rightarrow$ $\tfrac{h}{7},$ h = $7 \cdot \mathrm{sin}\,63^\circ$}{}
\end{pfloesung}
\begin{pfloesung}{}
\lz{102.}{$\approx 27{,}0$ m}{$45 \cdot \mathrm{tan}\,31^\circ$}{}
\end{pfloesung}
\begin{pfloesung}{}
\lz{103.}{57,9 m; 69,9 m}{112,8 · sin 30° $\Rightarrow$ 56,4 m; 56,4 m + 1,5 m $\Rightarrow$ 57,9 m}{}
\end{pfloesung}
\begin{pfloesung}{}
\lz{104.}{BD = 4,1 · sin 123° : sin 36° $\approx$ 5,85 cm}{BD : sin 123° $\Rightarrow$ 4,1 : sin 36° $\approx$ 5,85 cm}{}
\end{pfloesung}
\begin{pfloesung}{}
\lz{105.}{PB = 1,20 : tan 34,9° $\approx$ 1,72 m}{tan: 34,9° = 1,20 : PB}{}
\end{pfloesung}
\begin{pfloesung}{}
\lz{106.}{AC = 12 · sin 55° : sin 5° $\approx$ 112,8 m}{AC : sin 55° $\Rightarrow$ 12 : sin 5°}{}
\end{pfloesung}
\begin{pfloesung}{}
\lz{107.}{AB = 1500 · sin 55° $\approx$ 1228,7 m $\approx$ 1229 m}{sin: 55° = AB : 1500}{}
\end{pfloesung}
\begin{pfloesung}{}
\lz{108.}{AD = 1500 · cos 55° $\approx$ 860,4 m}{cos: 55° = AD : 1500}{}
\end{pfloesung}
\begin{pfloesung}{}
\lz{109.}{a = 7,4 : sin 52° $\approx$ 9,4 m}{sin: 52° $\approx$ 0,788}{}
\end{pfloesung}
\begin{pfloesung}{}
\lz{110.}{y = 160 · sin 141° : sin 34° $\approx$ 180,1 cm}{dritter Winkel: 34° gegenüber 160 cm}{}
\end{pfloesung}
\begin{pfloesung}{}
\lz{111.}{2800 · sin 60° : sin 50° + 2500 $\approx$ 5665 m ($\approx$ 5,67 km)}{2800 · sin 60° : sin 50° $\Rightarrow$ $\approx$ 3165 m}{}
\end{pfloesung}
\begin{pfloesung}{}
\lz{112.}{BD = √553 719 · tan 34° $\approx$ 502 m}{36°; Unterschied: 70° $-$ 36° $\Rightarrow$ 34°; AD · tan 34°}{}
\end{pfloesung}
\begin{pfloesung}{}
\lz{113.}{BC = 384 · sin 38° : sin 34° $\approx$ 422,8 m}{Unterschied: 180° $-$ 38° $-$ 108° $\Rightarrow$ 34°}{}
\end{pfloesung}
\begin{pfloesung}{}
\lz{114.}{BF = 131,5 · sin 45° = 65,75√2 $\approx$ 93,0 cm = AF; BC = 65,75√2 : cos 65° $\approx$ 220,0 cm}{sin: 45° = √$\tfrac{2}{2}$ $\approx$ 0,7071; cos: 65° $\approx$ 0,4226}{}
\end{pfloesung}
\begin{pfloesung}{}
\lz{115.}{DP = 11 · sin 115° : sin 49° $-$ 10 $\approx$ 3,2 m}{Unterschied: 180° $-$ 115° $-$ 16° $\Rightarrow$ 49°; DE : sin 115° $\Rightarrow$ 11 : sin 49° $\approx$ 13,2 m}{}
\end{pfloesung}
\begin{pfloesung}{}
\lz{116.}{BC = 2004 · sin 60° : sin 74° $\approx$ 1805,5 m; Aussage falsch (DC = 2004 m)}{Unterschied: 109° $-$ 35° $\Rightarrow$ 74°; BC : sin 60° $\Rightarrow$ 2004 : sin 74°}{}
\end{pfloesung}
\begin{pfloesung}{}
\lz{117.}{AD = 9 · sin 60° : sin 80° $\approx$ 7,9 m}{Unterschied: 90° $-$ 64° $\Rightarrow$ 26°; Unterschied: 86° $-$ 26° $\Rightarrow$ 60°; AD : sin 60° $\Rightarrow$ 9 : sin 80°}{}
\end{pfloesung}
\begin{pfloesung}{}
\lz{118.}{BC = 4,5 · sin 62,8° : sin 34,9° $\approx$ 7,00 m; 8 · (4,5 + BC) $\approx$ 92 m²}{BC : sin 62,8° $\Rightarrow$ 4,5 : sin 34,9° $\approx$ 7,00 m; Rechtecke: 36 m² und 56 m²}{}
\end{pfloesung}
\begin{pfloesung}{}
\lz{119.}{Sinussatz}{}{}
\end{pfloesung}
\begin{pfloesung}{}
\lz{120.}{$\approx 5{,}9$ cm}{$\tfrac{b}{\mathrm{sin}\,51^\circ}$ $\Rightarrow$ $\tfrac{7}{\mathrm{sin}\,68^\circ},$ b = $7 \cdot \mathrm{sin}\,51^\circ : \mathrm{sin}\,68^\circ$}{}
\end{pfloesung}
\begin{pfloesung}{}
\lz{121.}{kürzer als $1{,}7$ km, ja}{$B_1B_2$ = $1{,}8 \cdot \mathrm{sin}\,58^\circ : \mathrm{sin}\,67^\circ$ $\approx$ 1{,}66 km}{}
\end{pfloesung}
\begin{pfloesung}{}
\lz{122.}{57,9 m; 69,9 m}{112,8 · sin 30° $\Rightarrow$ 56,4 m; 56,4 m + 1,5 m $\Rightarrow$ 57,9 m}{}
\end{pfloesung}
\begin{pfloesung}{}
\lz{123.}{BD = 4,1 · sin 123° : sin 36° $\approx$ 5,85 cm}{BD : sin 123° $\Rightarrow$ 4,1 : sin 36° $\approx$ 5,85 cm}{}
\end{pfloesung}
\begin{pfloesung}{}
\lz{124.}{PB = 1,20 : tan 34,9° $\approx$ 1,72 m}{tan: 34,9° = 1,20 : PB}{}
\end{pfloesung}
\begin{pfloesung}{}
\lz{125.}{AC = 12 · sin 55° : sin 5° $\approx$ 112,8 m}{AC : sin 55° $\Rightarrow$ 12 : sin 5°}{}
\end{pfloesung}
\begin{pfloesung}{}
\lz{126.}{AB = 1500 · sin 55° $\approx$ 1228,7 m $\approx$ 1229 m}{sin: 55° = AB : 1500}{}
\end{pfloesung}
\begin{pfloesung}{}
\lz{127.}{AD = 1500 · cos 55° $\approx$ 860,4 m}{cos: 55° = AD : 1500}{}
\end{pfloesung}
\begin{pfloesung}{}
\lz{128.}{a = 7,4 : sin 52° $\approx$ 9,4 m}{sin: 52° $\approx$ 0,788}{}
\end{pfloesung}
\begin{pfloesung}{}
\lz{129.}{y = 160 · sin 141° : sin 34° $\approx$ 180,1 cm}{dritter Winkel: 34° gegenüber 160 cm}{}
\end{pfloesung}
\begin{pfloesung}{}
\lz{130.}{2800 · sin 60° : sin 50° + 2500 $\approx$ 5665 m ($\approx$ 5,67 km)}{2800 · sin 60° : sin 50° $\Rightarrow$ $\approx$ 3165 m}{}
\end{pfloesung}
\begin{pfloesung}{}
\lz{131.}{BD = √553 719 · tan 34° $\approx$ 502 m}{36°; Unterschied: 70° $-$ 36° $\Rightarrow$ 34°; AD · tan 34°}{}
\end{pfloesung}
\begin{pfloesung}{}
\lz{132.}{BC = 384 · sin 38° : sin 34° $\approx$ 422,8 m}{Unterschied: 180° $-$ 38° $-$ 108° $\Rightarrow$ 34°}{}
\end{pfloesung}
\begin{pfloesung}{}
\lz{133.}{BF = 131,5 · sin 45° = 65,75√2 $\approx$ 93,0 cm = AF; BC = 65,75√2 : cos 65° $\approx$ 220,0 cm}{sin: 45° = √$\tfrac{2}{2}$ $\approx$ 0,7071; cos: 65° $\approx$ 0,4226}{}
\end{pfloesung}
\begin{pfloesung}{}
\lz{134.}{DP = 11 · sin 115° : sin 49° $-$ 10 $\approx$ 3,2 m}{Unterschied: 180° $-$ 115° $-$ 16° $\Rightarrow$ 49°; DE : sin 115° $\Rightarrow$ 11 : sin 49° $\approx$ 13,2 m}{}
\end{pfloesung}
\begin{pfloesung}{}
\lz{135.}{BC = 2004 · sin 60° : sin 74° $\approx$ 1805,5 m; Aussage falsch (DC = 2004 m)}{Unterschied: 109° $-$ 35° $\Rightarrow$ 74°; BC : sin 60° $\Rightarrow$ 2004 : sin 74°}{}
\end{pfloesung}
\begin{pfloesung}{}
\lz{136.}{AD = 9 · sin 60° : sin 80° $\approx$ 7,9 m}{Unterschied: 90° $-$ 64° $\Rightarrow$ 26°; Unterschied: 86° $-$ 26° $\Rightarrow$ 60°; AD : sin 60° $\Rightarrow$ 9 : sin 80°}{}
\end{pfloesung}
\begin{pfloesung}{}
\lz{137.}{BC = 4,5 · sin 62,8° : sin 34,9° $\approx$ 7,00 m; 8 · (4,5 + BC) $\approx$ 92 m²}{BC : sin 62,8° $\Rightarrow$ 4,5 : sin 34,9° $\approx$ 7,00 m; Rechtecke: 36 m² und 56 m²}{}
\end{pfloesung}
\begin{pfloesung}{}
\lz{138.}{1~Pythagoras, 2~Winkelfunktion, 3~Seite berechnen, 4~Winkel, 5~Pythagoras, 6~Winkelfunktion}{}{}
\end{pfloesung}
\begin{pfloesung}{}
\lz{139.}{gleichschenklig}{}{}
\end{pfloesung}
\begin{pfloesung}{}
\lz{140.}{Scheitelwinkel}{}{}
\end{pfloesung}
\begin{pfloesung}{}
\lz{141.}{Dreieck ADC}{}{}
\end{pfloesung}
\begin{pfloesung}{}
\lz{142.}{$38^\circ$}{Scheitelwinkel sind gleich groß}{}
\end{pfloesung}
\begin{pfloesung}{}
\lz{143.}{4 · 80° = 320° $\neq$ 360° (Winkelsumme im Viereck)}{}{}
\end{pfloesung}
\begin{pfloesung}{}
\lz{144.}{z. B. $\beta$ = 60°, $\gamma$ = 57° ($\beta$ + $\gamma$ = 117°)}{}{}
\end{pfloesung}
\begin{pfloesung}{}
\lz{145.}{$\gamma$ = 180° $-$ 110° $-$ 25° = 45°}{}{}
\end{pfloesung}
\begin{pfloesung}{}
\lz{146.}{$116^\circ$}{$116^\circ, \gamma$ $\Rightarrow$ $64^\circ, \delta$}{}
\end{pfloesung}
\begin{pfloesung}{}
\lz{147.}{$\beta$ = 105°}{Unterschied: 180° $-$ 75°}{}
\end{pfloesung}
\begin{pfloesung}{}
\lz{148.}{$60^\circ$}{$180^\circ : 3$}{}
\end{pfloesung}
\begin{pfloesung}{}
\lz{149.}{$\alpha$ = 40°}{}{}
\end{pfloesung}
\begin{pfloesung}{}
\lz{150.}{$\gamma$ = 82,3°}{}{}
\end{pfloesung}
\begin{pfloesung}{}
\lz{151.}{$\delta$ = 21°}{}{}
\end{pfloesung}
\begin{pfloesung}{}
\lz{152.}{$\gamma$ = 124°}{Unterschied: 180° $-$ 56°}{}
\end{pfloesung}
\begin{pfloesung}{}
\lz{153.}{$\sphericalangle$ACB = 60° $\Rightarrow$ $\sphericalangle$ACD = 120° $\Rightarrow$ $\delta$ = 180° $-$ 120° $-$ 5° = 55°}{}{}
\end{pfloesung}
\begin{pfloesung}{}
\lz{154.}{$\gamma_1$ = cos$^{-1}$($\tfrac{74}{141}$) $\approx$ 58,3°}{Unterschied: 180° $-$ 90° $-$ 31,7°}{}
\end{pfloesung}
\begin{pfloesung}{}
\lz{155.}{$2$ Achsen: die beiden Mittellinien}{}{}
\end{pfloesung}
\begin{pfloesung}{}
\lz{156.}{2}{}{}
\end{pfloesung}
\begin{pfloesung}{}
\lz{157.}{$3$\,cm}{6 : 2 $\Rightarrow$ 3}{}
\end{pfloesung}
\begin{pfloesung}{}
\lz{158.}{die senkrechte Diagonale durch die obere und die untere Spitze}{}{}
\end{pfloesung}
\begin{pfloesung}{}
\lz{159.}{$1$}{Drachenviereck}{}
\end{pfloesung}
\begin{pfloesung}{}
\lz{160.}{Parallelogramm}{}{}
\end{pfloesung}
\begin{pfloesung}{}
\lz{161.}{$3$}{}{}
\end{pfloesung}
\begin{pfloesung}{}
\lz{162.}{$4$, $2$, $2$, $1$}{Quadrat, Rechteck, Raute, Drachenviereck}{}
\end{pfloesung}
\begin{pfloesung}{}
\lz{163.}{4}{}{}
\end{pfloesung}
\begin{pfloesung}{}
\lz{164.}{1}{Mittelsenkrechte der parallelen Seiten}{}
\end{pfloesung}
\begin{pfloesung}{}
\lz{165.}{1 Symmetrieachse (BD); BF = 80 $-$ 25 = 55 cm; AB = 5√146 $\approx$ 60,4 cm}{25 cm; 25² + 55² $\Rightarrow$ 3650 cm²}{}
\end{pfloesung}
\begin{pfloesung}{}
\lz{166.}{$40$ cm}{Hypotenuse ist KM; 2{,}6 cm; Kathete gesucht: $KL^2$ $\Rightarrow$ 42{,}25 - 10{,}89, KL = 5{,}6 dm; $41 cm, LM^2$ $\Rightarrow$ 1\,681 - 81, LM; gegenüber L; a; b}{}
\end{pfloesung}
\begin{pfloesung}{}
\lz{167.}{$\approx 46{,}6^\circ$}{$25{,}3^\circ$; Tangens mit: $60 cm, \alpha$ $\approx$ $36{,}9^\circ$; Kosinus mit: $16 cm, \alpha$; 1; 2; 3}{}
\end{pfloesung}
\begin{pfloesung}{}
\lz{168.}{$24$ cm}{$\sqrt{36 + 4 + 81}$ $\Rightarrow$ 11 cm; Überstand: 5 cm, h = $\sqrt{13^2 - 5^2}$ = 12 cm; Radius: 7 cm, h = $\sqrt{25^2 - 7^2}$; 1; 3}{}
\end{pfloesung}
\begin{pfloesung}{}
\lz{169.}{$\approx 6{,}2$ cm}{rechtwinklig, Sinus: $8 \cdot \mathrm{sin}\,40^\circ$ $\Rightarrow$ $\approx$ 5{,}1 cm; Sinussatz: $\tfrac{6 \cdot \mathrm{sin}\,75^\circ}{\mathrm{sin}\,50^\circ}$ $\Rightarrow$ $\approx$ 7{,}6 cm; Kosinussatz: $\sqrt{39}$; 1; 3}{}
\end{pfloesung}
\begin{pfloesung}{}
\lz{170.}{Im gleichseitigen Dreieck sind alle Winkel gleich: 180° : 3 = 60°; ein 90°-Winkel ist nicht möglich.}{}{}
\end{pfloesung}
\begin{pfloesung}{}
\lz{171.}{Quadrat, Rechteck, Raute}{}{}
\end{pfloesung}
\begin{pfloesung}{}
\lz{172.}{gibt es zwei Paare gleich langer Nachbarseiten}{}{}
\end{pfloesung}
\begin{pfloesung}{}
\lz{173.}{gegenüberliegende Winkel gleich groß}{}{}
\end{pfloesung}
\begin{pfloesung}{}
\lz{174.}{gibt es ein Paar paralleler Seiten}{Definition Trapez}{}
\end{pfloesung}
\begin{pfloesung}{}
\lz{175.}{$46^\circ$}{Unterschied: $180^\circ - 2 \cdot 67^\circ$}{}
\end{pfloesung}
\begin{pfloesung}{}
\lz{176.}{BC = 5 cm}{gleiche Winkel bei A und B}{}
\end{pfloesung}
\begin{pfloesung}{}
\lz{177.}{BC = 6 cm ($\alpha$ = $\beta$ $\Rightarrow$ gleichschenklig); $\gamma$ = 40°}{Unterschied: 180° $-$ 2 · 70°}{}
\end{pfloesung}
\begin{pfloesung}{}
\lz{178.}{Winkel bei A und B im Dreieck ABF je 45° $\Rightarrow$ gleichschenklig mit Basis AB (BF = AF)}{Winkelsumme: 180° $-$ 90° $-$ 45° $\Rightarrow$ 45°}{}
\end{pfloesung}
\begin{pfloesung}{}
\lz{179.}{Seitenhalbierende}{}{}
\end{pfloesung}
\begin{pfloesung}{}
\lz{180.}{Ja}{6² + 8² = 100 = 10²}{}
\end{pfloesung}
\begin{pfloesung}{}
\lz{181.}{AF = DF = 25 cm $\Rightarrow$ Dreieck AFD gleichschenklig-rechtwinklig $\Rightarrow$ $\angle$ADF = 45°, ebenso $\angle$FDC = 45° $\Rightarrow$ $\gamma$ = 90° (oder Thales: D auf Kreis um F mit r = 25 cm)}{Basiswinkel: 45°}{}
\end{pfloesung}
\begin{pfloesung}{\textbf{Prüfstein}}
\lz{a)}{h = 3√95 $\approx$ 29,2 cm}{Unterschied: 32² $-$ 13² $\Rightarrow$ 855}{}
\lz{b)}{$\varepsilon$ = cos$^{-1}$($\tfrac{13}{32}$) $\approx$ 66,0°}{Anteil: 13 : 32 $\Rightarrow$ $\tfrac{13}{32}$ $\approx$ 0,406}{}
\lz{c)}{AC = 3√95 : sin 22° $\approx$ 78,1 cm}{sin: 22°; 3√95 : tan 22° $\Rightarrow$ $\approx$ 72,4 cm}{}
\end{pfloesung}
\end{document}